\section{Construction requirements and the next questions}\label{sec:roadmap}

The two groups of results constrain different objects. The cover
theorems concern the coefficient algebra of a homogeneous check.
The local theorem concerns the span of localized words in the code
itself. The theorems do not supply either premise from the other.
Their common use is to
separate a promising construction from a family whose capacity is
already controlled.

\begin{table}[htbp]
\centering\small
\begin{tabular}{@{}>{\raggedright\arraybackslash}p{4.0cm}>{\raggedright\arraybackslash}p{4.5cm}>{\raggedright\arraybackslash}p{6.0cm}@{}}
\toprule
Hypothesis required of a candidate & Consequence here & What an above-half candidate must change \\
\midrule
Fixed separated template, specified pro-$p$ chain, homogeneous fixed widths
& Half rate over every characteristic-$p$ field
& The representation, homogeneity, fixed-width condition, or code class. \\[4pt]
Same template and chain, growing widths with tight simple-factor mass
& Half rate; excess $\delta$ forces at least $2\delta$ large-factor mass
& A positive fraction of dimension in arbitrarily large absolute simple factors. \\[4pt]
Independent-chord check, nonzero scalar edges, or invertible common-flag blocks
& Rate tends to zero
& Edge nonvanishing, the common flag, or the specified check family. \\[4pt]
Characteristic two, odd girth greater than $2R+1$, radius-$R$ local generation
& Weighted half rate
& A quantitatively nonlocal code subspace, or the graph premise. \\
\bottomrule
\end{tabular}
\caption{These are conditional tests, not automatic properties of tape computations.}
\label{tab:requirements}
\end{table}
\FloatBarrier

\subsection{A focused constructive target}
The first target is a fixed separated voltage template and an actual
chain of the stated type, with homogeneous coefficients whose
geometric simple-factor dimensions grow. An above-half sequence
would have to meet the mass lower bound of \cref{cor:factor-excess}.
Large factors alone are not evidence of high rate: the exact
inflation control in \cref{prop:large-factor} has them and still
has a separate all-code half bound. A candidate should therefore be
tested both for large factors and for an independent coding obstruction.

If the search instead uses characteristic-two codes on graphs of
large odd girth, \cref{cor:nonlocal} gives a second test. Local
recovery equations must leave enough independent global codewords
outside the local span. A short description of the equations or a
sparse check matrix is insufficient; the needed object is the
actual kernel or subcode, with its dimension proved.

\subsection{The corresponding impossibility targets}
\begin{question}[Unrestricted block width]\label{prob:all-width}
Under the fixed template, separation, embedding, and chain hypotheses
of \cref{thm:fixed-width}, does half rate hold for every sequence of
homogeneous characteristic-$p$ checks, without a condition on their
simple-factor dimensions?
\end{question}
An affirmative answer would, by \cref{prop:inflation}, also control
arbitrary sheet-dependent linear checks on those same quotient
relations. Conversely, such a sheet-dependent above-half check
would inflate to a homogeneous counterexample with growing width.
This is an equivalence within the specified cover model; it does
not extend the template or group hypotheses.

\begin{question}[A native source of the needed structure]
\label{prob:native}
Can actual shallow-cone fibres of a putative linear-time transposer
be placed, with a fully charged cost, in one of the code classes
controlled here or in an all-alphabet extension of those classes?
\end{question}
The native paired-distance theorem of~\cite{BrendeDirected} alone
does not give a fixed voltage template, homogeneous recovery,
linearity of the fibre, an internal flag, or locally supported
generators. These are the main missing implications. Even a
successful coding step must be combined with a compatible cheap
depth reduction and, for adaptive machines, a paid description of
the input-dependent trace or relation.
